\chapter{Contexte général du projet}
\label{ch:contexte}

\section{Organisme d'accueil}

IAT Academy est une académie de formation professionnelle spécialisée dans
les métiers de l'aviation, de l'accueil et du tourisme, préparant en
particulier aux filières Hôtesse de l'air et Steward, Agent d'escale,
Accompagnateur touristique, et Réception hôtelière. La formation se
déroule sur un cycle professionnel de deux années, structuré autour d'un
équilibre volontairement assumé entre théorie et pratique : cours en
présentiel dispensés par des formateurs issus directement du secteur
(compagnies aériennes, groupes hôteliers, agences de voyage), stage en
milieu réel intégré au cursus plutôt que relégué à une simple option de
fin de parcours, et, désormais, un espace d'apprentissage numérique
prolongeant la salle de classe au-delà des heures de cours.

Cette double identité --- académie de formation aux métiers du transport
aérien et du tourisme, mais aussi organisme dont l'indicateur de réussite
premier reste l'insertion professionnelle effective de ses diplômés ---
irrigue l'ensemble des choix fonctionnels du projet. Elle explique
pourquoi le référentiel de compétences se traduit directement dans
l'organisation du contenu pédagogique (unités de formation alignées sur
les compétences métier attendues par les employeurs du secteur), pourquoi
le stage occupe une place aussi centrale dans le modèle de progression
(voir chapitre~\ref{ch:analyse-conception}), et pourquoi la plateforme
intègre jusqu'à un mécanisme de certification vérifiable publiquement et
de partage professionnel sur les réseaux sociaux (voir
chapitre~\ref{ch:realisation}) --- un choix qui, sur un plan strictement
académique, pourrait sembler périphérique, mais qui prend tout son sens
dès lors que l'on considère la mission réelle de l'académie : préparer ses
apprenants non pas seulement à réussir des examens, mais à être recrutés.
L'identité visuelle de la plateforme reprend d'ailleurs ce parti pris
métier jusque dans son vocabulaire d'interface, construit autour de la
métaphore de la carte d'embarquement --- un choix esthétique qui ancre
délibérément l'expérience numérique dans l'univers professionnel visé par
les apprenants.

\section{Contexte et problématique}

Avant ce projet, le suivi pédagogique d'IAT Academy reposait sur des
supports et des outils non centralisés : le contenu de cours, l'évaluation
des apprenants et le suivi de leur progression n'étaient réunis dans aucun
espace numérique unique. Cette situation, courante dans de nombreux
établissements de formation professionnelle de taille intermédiaire,
posait en pratique plusieurs difficultés concrètes que ce projet s'est
attaché à résoudre une à une.

La première difficulté tenait à l'absence de toute vue d'ensemble, en
temps réel, de la progression d'un apprenant à travers les deux années de
son cursus. Sans outil centralisé, savoir précisément où en est un
apprenant donné --- quels modules il a complétés, quelles évaluations lui
restent à passer, si son stage a été validé --- exigeait de recouper
manuellement plusieurs sources d'information, un exercice difficilement
soutenable à l'échelle de plusieurs promotions. La seconde difficulté
concernait l'évaluation elle-même : les quiz, devoirs et réponses libres
des apprenants n'étaient ni centralisés, ni tracés de façon homogène,
rendant la correction et le suivi des notes dépendants de pratiques
individuelles plutôt que d'un processus outillé et cohérent.

La troisième difficulté, sans doute la plus structurante pour la
conception technique du projet, était l'absence de tout mécanisme fiable
de validation d'étape : rien, dans l'organisation existante, n'empêchait
techniquement un apprenant d'accéder à du contenu avancé sans avoir
réellement validé les prérequis correspondants --- une lacune d'autant
plus problématique que certaines étapes du cursus d'IAT Academy, comme la
validation du stage en entreprise ou de la soutenance finale, ne peuvent
pas se réduire à un critère automatique et exigent le jugement explicite
du Directeur pédagogique. Enfin, la délivrance et la vérification des
certificats de réussite reposaient sur un processus entièrement manuel,
sans preuve numérique qu'un tiers extérieur --- un employeur, un
partenaire du secteur --- puisse vérifier de façon autonome.

Cette analyse conduit à formuler la problématique centrale de ce projet de
la manière suivante : \emph{comment structurer, sécuriser et automatiser
le parcours pédagogique complet d'un apprenant --- de l'inscription à la
délivrance du certificat --- tout en respectant des règles métier qui ne
se réduisent jamais à de simples critères automatiques, puisque certaines
étapes-clés exigent explicitement l'implication du Directeur
pédagogique ?} C'est cette tension entre automatisation et jugement
humain, plus que la seule numérisation d'un contenu de cours, qui a guidé
l'ensemble des choix de conception présentés au
chapitre~\ref{ch:analyse-conception}.

\section{Objectifs du projet}

Pour répondre à cette problématique, six objectifs concrets ont structuré
le travail réalisé durant ce stage.

Le premier objectif consistait à concevoir un modèle de données
représentant fidèlement la structure pédagogique réelle du cursus ---
Formation, Année, Unité de Formation, Module, Leçon, puis Bloc de contenu
--- sans pour autant complexifier inutilement le schéma relationnel sous-jacent ;
ce choix de modélisation, en apparence anodin, s'est révélé être l'une des
décisions de conception les plus structurantes du projet, comme le détaille
la section~\ref{sec:analyse-conception} du chapitre suivant. Le deuxième
objectif visait à mettre en place un système d'évaluation complet, capable
de gérer plusieurs types de quiz selon leur portée --- section, module, fin
d'unité de formation, fin d'année --- ainsi que plusieurs types de
questions, avec à la fois une correction automatique pour les questions à
choix et une correction manuelle pour les réponses libres, ces deux modes
devant in fine converger vers un seul et même score final cohérent.

Le troisième objectif était de garantir un déblocage séquentiel du contenu
pédagogique fondé sur des règles vérifiées côté serveur, et non pas
seulement suggérées côté interface --- une exigence de sécurité autant que
de fiabilité pédagogique. Le quatrième objectif consistait à fournir à
l'équipe pédagogique d'IAT Academy un espace d'administration complet,
lui permettant de produire le contenu de cours, de corriger les
évaluations, de suivre individuellement chaque apprenant et de réaliser les
validations manuelles qui, comme évoqué plus haut, ne peuvent être
déléguées à un algorithme. Le cinquième objectif visait à automatiser la
délivrance de certificats numériquement vérifiables, tout en outillant le
suivi du stage jusqu'à sa signature par un tuteur externe à l'académie.
Enfin, le sixième objectif portait sur la sécurisation transversale de
l'ensemble de la plateforme --- authentification, séparation stricte des
rôles, protection contre les abus des fonctionnalités exposées à des
services tiers --- et sur son déploiement de façon reproductible, condition
nécessaire à une exploitation réelle au-delà du cadre du stage.

\section{Méthodologie de travail et planning}

Le développement du projet a suivi une démarche itérative, par
fonctionnalité plutôt que par couche technique : chaque bloc fonctionnel
--- authentification, catalogue de cours, système de quiz, gestion du
stage, certification, puis les fonctionnalités transverses de
communication et de sécurité --- a été conçu, implémenté et testé de bout
en bout (du modèle de données jusqu'à l'interface) avant de passer au
suivant, plutôt que de développer séparément l'intégralité du backend puis
l'intégralité du frontend. Cette approche a permis de valider très tôt,
sur des cas concrets, que le modèle de données retenu tenait effectivement
la charge des règles métier réelles --- en particulier la logique de
progression par unité de formation --- avant d'investir du temps de
développement sur les fonctionnalités qui en dépendent.

Des points d'avancement réguliers avec l'encadrant professionnel au sein
d'IAT Academy ont permis de confronter les choix techniques aux
contraintes réelles de l'académie, notamment sur les points où une règle
métier n'était pas immédiatement évidente à traduire en logique
applicative --- la distinction entre les unités de formation exigeant une
validation manuelle du Directeur et celles n'en exigeant pas en est
l'exemple le plus révélateur (voir section~\ref{sec:progression-uf}).

Le temps du stage s'est réparti, dans ses grandes masses, entre une phase
initiale de cadrage et d'analyse des besoins auprès de l'équipe
pédagogique, une phase de conception du modèle de données et de
l'architecture --- la plus déterminante pour la suite, dans la mesure où
une erreur de modélisation à ce stade se serait répercutée sur l'ensemble
du développement ---, une phase de développement à proprement parler,
volontairement la plus longue, menée fonctionnalité par fonctionnalité
comme décrit ci-dessus, et enfin une phase de consolidation regroupant les
tests, la correction des anomalies détectées et la rédaction du présent
rapport. Le tableau~\ref{tab:planning} synthétise cette répartition ; il
peut être complété par un diagramme de Gantt détaillé si le format attendu
par l'établissement l'exige.

\begin{table}[H]
  \centering
  \caption{Répartition indicative des grandes phases du projet}
  \label{tab:planning}
  \begin{tabular}{@{}p{5cm}p{8cm}@{}}
    \toprule
    Phase & Contenu principal \\
    \midrule
    Cadrage et analyse des besoins & Immersion dans le fonctionnement d'IAT Academy, recueil des règles métier, définition du périmètre. \\
    Conception & Modélisation des entités, choix d'architecture, définition des mécanismes de sécurité. \\
    Développement & Implémentation backend et frontend, fonctionnalité par fonctionnalité. \\
    Consolidation & Rédaction et exécution des tests, correction des anomalies, rédaction du rapport. \\
    \bottomrule
  \end{tabular}
\end{table}

\section{Conclusion}

Ce chapitre a présenté le contexte métier d'IAT Academy, la problématique
qui a motivé ce projet --- automatiser un parcours pédagogique complexe
sans sacrifier les validations qui exigent un jugement humain --- ainsi
que les objectifs et la méthodologie qui en découlent. Le
chapitre~\ref{ch:analyse-conception} détaille maintenant l'analyse
fonctionnelle et les choix de conception qui répondent concrètement à ces
objectifs.
